The execution records below document data handling and the timing of analysis decisions. They supplement the scientific procedures in the main Methods; they do not alter the adaptive or exploratory status of any comparison.

\subsection*{FER2013 acquisition and FER+ alignment}
We acquired the CSV on 20 September 2026 from Kaggle mirror \texttt{deadskull7/fer2013}, version 1, not from an official maintainer distribution. The join follows the official FER+ conversion script's row indexing, not its output filenames, which can be blank or duplicated. A separate raw-source audit checks all 35,887 original emotion IDs, pixel byte hashes and partition positions against the retained manifest. We reorder named vote columns to the fixed model class order. Online Resource 1 contains acquisition and split manifests, the original annotation records, the row adapter and the raw-source audit.

\subsection*{External protocol and integrity records}
We froze the local protocol on 11 October 2026 at 18:17:50 UTC, before computing external model outcomes. It records implementation, annotation, detector and checkpoint hashes. Internal outcomes and external annotation summaries were already known. The pre-outcome amendment corrects the disgust/surprise response mapping using primary-source descriptors and accuracy records; the original freeze is retained. Each RAVDESS actor archive passed the publisher's MD5 checksum, and each extracted clip retains a SHA-256 hash. Online Resource 8 includes the freeze, amendment, integrity manifests, annotation adapter and retained frame/clip predictions.

\subsection*{Post hoc diagnostic addendum}
A dated addendum on 12 October 2026 specified these descriptive outputs after the internal and external results were known. The addendum covers category frequencies, vote-weighted reference/prediction matrices and temporal summaries. Its status remains post hoc: no new training, fitting, frame-peak selection or confirmatory test is introduced.

\subsection*{Final follow-up precision and execution repair}
The 20 October follow-up retains its original scientific protocol and freeze. Its target adapter computes normalized supported proportions in float64 before float32 target and weight construction. Before permutation, comparing reconstructed seven-class Soft coefficients with the retained historical float32 raw-vote implementation identifies differences on 349 training rows, with maximum absolute difference $2.98\times10^{-8}$ and mean $5.29\times10^{-11}$ across coefficients. The annotations are unchanged. These arithmetic differences preclude claiming bitwise-identical loss coefficients across historical and new implementations; new eight-output conditions share the same adapter.

After Permutation seed 17 retained epoch 1, the next epoch failed when Windows could not allocate a worker's CPU tensor. Amendment 001 replaced spawned workers with sequential loading before new calibration/test inference. Its adapter preserves the persistent-loader sample order, stateless flips and optimizer/model RNG states, resuming after the committed epoch without repeating its updates. The source, amendment freeze and three-epoch comparison against an actual two-worker loader are retained.

Amendment 002 added one input-prefetch thread, bounded to one additional batch, and one intra-op CPU thread, matching the former workers. Sampler and transforms are unchanged. Permutation seed 17 resumed after committed epoch 2; every run records its amendment timing. A second actual two-worker comparison verifies all sample positions over three synthetic epochs. Both amendments, adapters, freezes and tests accompany the release. They preceded new calibration/test inference and changed no target, image, seed, batch size, optimizer, schedule, budget or endpoint.

An asynchronous CUDA error later interrupted selection inference in epoch 12 of eight-output Uniform, seed 89. The unchanged adapter successfully replayed that partial epoch from the verified epoch-11 model, optimizer, scheduler and RNG checkpoint. A separate recovery record reports the interruption without assigning a hardware or software cause.
